\documentclass[12pt,a4paper]{amsart}
%\input xy
%\xyoption{all}
\usepackage[applemac]{inputenc}

\DeclareMathAlphabet{\mathpzc}{OT1}{pzc}{m}{it}
\usepackage{amsfonts,amsmath,amssymb,indentfirst,mathrsfs,amscd,url}
\usepackage[mathscr]{eucal} 
%\usepackage[active]{srcltx}




\setlength{\oddsidemargin}{0.4in}
\setlength{\evensidemargin}{0.4in}
\setlength{\textwidth}{5.5in}
\setlength{\textheight}{8.8in}
\setlength{\marginparwidth}{0.8in}
\addtolength{\headheight}{2.5pt}

\renewcommand{\thefootnote}{}

\begin{document}
\thispagestyle{empty}
\vspace*{-2cm}

 
\noindent
\begin{center}\Large\bfseries\textsf{Curso avanzado de Geometr�a\\
M\'aster en matem\'aticas,   curso 2013-14\\}
\end{center}
\hrulefill{}
 
\noindent
{\sffamily    
Profesor:} Marco Zambon. \\
%{\sffamily Despacho:} 01.17.AU.103\\
{\sffamily Tutor\'ias:} por cita previa, en el despacho 605 del departamento de matem\'aticas (modulo 17). \\
{\sffamily P\'agina web del curso:} \url{http://www.uam.es/marco.zambon/SS14.html}

\noindent\hrulefill{}  

\noindent
{\sffamily Horario:} Martes y Jueves, 11:30-13:00, en el aula 320 del departamento de matem\'aticas.    

%\noindent
%{\sffamily  Fecha de examenes:} \\ 
%ejercicios realizados en clase: a decidir\\
%examen final: 8 de enero de 2014 (miercoles, tarde)\\
%convocatoria extraordinaria:  20 de junio de 2014 (viernes, ma\~nana) 
\noindent\hrulefill{}  

\noindent
{\sffamily  Sistema de evaluaci\'on:} 
% (Ver la Guia Docente para m\'as detalles.)\\
La nota se determina a partir de:\\
entrega de ejercicios (70 por ciento),\\
breve examen oral  (30 por ciento). El examen tendr\'a lugar en la semana del 19 de mayo. Consistir\'a de la explicaci\'on de unos ejercicios entre los asignados a lo largo del semestre, y   de preguntas sobre una parte del temario que ser\'a concretada a lo largo del curso.

 \noindent\hrulefill{}  
  

 \noindent
{\sffamily Programa orientativo del curso:} \\

\noindent PARTE 1: GEOMETR�A SIMPL�CTICA\\

1A) Motivaci�n f�sica de la geometr�a simpl�ctica: la mec�nica cl�sica. Algebra lineal simpl�ctica: espacios vectoriales simpl�cticos y sus subespacios. Revisi�n: variedades, subvariedades, campos vectoriales. Variedades simpl�cticas. Obstrucciones topol�gicas. El teorema de Darboux. Subvariedades lagrangianas, coisotropicas. Formas normales: teoremas de Weinstein y de Gotay.\\

1B) Revisi�n: grupos de Lie. Acciones Hamiltonias y aplicaciones momento. Interpretaci�n en t�rminos de cohomolog�a equivariante (� la Atiyah-Bott). Reducci�n simpl�ctica. El teorema de convexidad de Atiyah-Guillemin-Sternberg.\\

\noindent PARTE 2: GEOMETR�A DE POISSON\\

2A)Variedades de Poisson: caracterizaci�n en t�rminos de hojas simpl�cticas, y en t�rminos de tensores de Poisson. Variedades de Poisson: caracterizaci�n en t�rminos de �lgebras de Poisson. Relaci�n con �lgebras de Lie.	Aplicaciones de Poisson y
subvariedades coisotropicas. Foliaci�nes singulares.\\

2B) Grupoides y algebroides de Lie. Ideas b�sicas sobre grupoides y algebroides de Lie associados a una variedad de Poisson. Ideas b�sicas sobre cuantizaci�n por deformaci�n.
 

\newpage
\noindent\hrulefill{}  

 \noindent
{\sffamily Bibliograf\'ia:} \\

\begin{itemize}
\item	Ana Cannas da Silva, ``Lectures on symplectic geometry", Springer Verlag. Disponible en\\ \url{http://www.math.ethz.ch/~acannas/Papers/lsg.pdf}
\item Eckhart Meinrenken,``Symplectic geometry", disponible en\\ \url{http://www.math.toronto.edu/mein/teaching/lectures.html}
\item Mich\`ele Audin, ``Torus Actions on Symplectic Manifolds'', Birkh\"auser Verlag, Progress in Mathematics, 2004.
\item Ana Cannas da Silva and Alan Weinstein``Geometric Models for Noncommutative Algebras", American Mathematical Society. Disponible en\\ \url{http://math.berkeley.edu/~alanw/}
\item Camille Laurent-Gengoux, Anne Pichereau, Pol Vanhaecke,``Poisson Structures", Grundlehren der mathematischen Wissenschaften. Disponible en la biblioteca
\item Izu Vaisman, ``Lectures on the geometry of Poisson manifolds", Birkh\"auser Verlag.
\end{itemize}

   
  

  
  
\end{document}
